\documentclass[11pt]{article}
\usepackage{amsmath,amssymb,geometry}
\geometry{margin=22mm}
\begin{document}

\section*{P52: full-fleet battery and charging design, explicit robust branch}
The source's nominal and robust equations cannot be imposed together with the
extracted opposite energy sign. The current branch follows physical energy
conservation and source (24)--(28), with independent upper consumption on every
leg. It is a conservative planning formulation.

\subsection*{Finite data and variables}
Routes are $r\in R$, their entire bus fleets $b\in B_r$, battery choices $k\in BAT$,
terminals $t\in T$ and charger choices $c\in FCS$. Every bus has an ordered complete
leg list $h=1,\ldots,H_b$, including its final return. Leg $h$ ends at $t_{bh}$
and has admissible charging cells $\Omega_{bh}$; its overlap duration in cell
$w$ is $a_{bhw}\in(0,\Delta]$ hours. No charging variable exists outside these
windows. The supplied worst consumption is
\[
 E_{bh}^{+}=\bar E_{bh}+\min(1,\Gamma_r)\widehat E_{bh}.
\]
Its value is already multiplied by the actual leg distance in the completed data.
This is the exact maximum of the single-coordinate uncertainty expression, used
conservatively on every leg, with a separate uncertainty budget per leg.

Binary $x_{rk}$ selects a battery; binary $y_{tc}$ selects a charger type;
$n_{tc}\in\mathbb Z_{\ge0}$ counts chargers. Binary $z_{bh}$ and $v_{bhw}$
mark charging events/cells. $p_{bhw}\ge0$ is the planned power in kW,
$A_{bh},D_{bh}\ge0$ arrival/departure energy in kWh, and $P_t^{on},P_t^{off}\ge0$
are the two demand-charge peaks. Define
$C_r=\sum_k CAP_kx_{rk}$, $M_t=\sum_{r:t\in T_r}|B_r|$ and
$P^{max}=\max_c CP_c$.

\subsection*{All constraints}
\begin{align}
\sum_kx_{rk}&=1,&
\sum_c y_{tc}&\le1,&
0\le n_{tc}&\le M_t y_{tc},\\
A_{b1}&=SOC_{max}C_r-E_{b1}^{+},&
A_{bh}&=D_{b,h-1}-E_{bh}^{+}\quad(h>1),\\
SOC_{min}C_r&\le A_{bh}\le SOC_{max}C_r,&
SOC_{min}C_r&\le D_{bh}\le SOC_{max}C_r,\\
D_{bh}&=A_{bh}+\sum_{w\in\Omega_{bh}}a_{bhw}p_{bhw},\\
0\le p_{bhw}&\le P^{max}v_{bhw},&
p_{bhw}&\le\sum_cCP_c y_{t_{bh},c},&
v_{bhw}&\le\sum_cy_{t_{bh},c},\\
v_{bhw}&\le z_{bh},&
z_{bh}&\le\sum_{w\in\Omega_{bh}}v_{bhw},&
z_{bh}&\le\sum_c y_{t_{bh},c}.
\end{align}
For each terminal and cell, let
$I_{tw}=\{(b,h):t_{bh}=t,\ w\in\Omega_{bh}\}$ and let $g(w)$ be its on/off
tariff category. Impose
\begin{align}
\sum_{(b,h)\in I_{tw}}v_{bhw}&\le\sum_c n_{tc},\\
\sum_{(b,h)\in I_{tw}}p_{bhw}&\le\sum_c CP_cn_{tc},\\
P_t^{g(w)}&\ge\Delta^{-1}\sum_{(b,h)\in I_{tw}}a_{bhw}p_{bhw}.
\end{align}
Counting a partial-cell charging event for the entire occupancy cell is
conservative; duration-weighted energy and demand averages use actual overlap.

\subsection*{Objective and units}
Write $a(L)=d/[1-(1+d)^{-L}]$ for annualization at discount $d>0$.
With supplied prices in dollars per kWh and demand rates in dollars per
kW-month, minimize
\begin{align}
&\sum_{r,k}a(L_k^B)C_k^B|B_r|x_{rk}
+\sum_{t,c}[a(L^{station})C^{station}+M^{station}]y_{tc}\\
&+\sum_{t,c}[a(L_c^C)C_c^C+M_c^C]n_{tc}
+365\sum_{b,h,w}c_w a_{bhw}p_{bhw}\\
&+12\sum_t[d^{on}P_t^{on}+d^{off}P_t^{off}].
\end{align}
Thus the monthly demand rate absorbs the source's daily-to-month conversion;
the monthly rate is applied once per month. Battery procurement
is fleet size times battery cost. Charger maintenance, if present in the completed
input, is explicitly included. The given full initial SOC is exogenous:
the energy bill covers daytime opportunity charging.

\subsection*{Robust operational meaning}
The optimization uses the upper-consumption trajectory. Under any smaller
realized nonnegative leg consumptions, charge no more than the planned
cell power, curtailing at $SOC_{max}C_r$. If actual arrival energy is $a$ and
the planned cell energy is $q$, the feedback update is
$a^+=\min\{SOC_{max}C_r,a+q\}$. Induction gives actual energy at least the
upper-consumption trajectory and at most the SOC ceiling; resource use and
cost do not increase. This curtailment policy provides a feasible operating policy and a conservative cost upper bound.

\end{document}
